\subsection{Direct similarities in a plane.}

Recall (§ 6, no. 5) that a direct similarity of E is an automorphism u of the vector space E such that \(\Phi(u(x), u(y)) = (\det u)\Phi(x, y)\) for all \(x, y\) in E.

PROPOSITION 1. --- Let A(Φ) be the subalgebra of \(\mathfrak{L}_{\mathbf{A}}(\mathbf{E})\) generated by the direct similarities of E.

\begin{enumerate}
\def\labelenumi{\alph{enumi})}
\item
  Direct similarities are the invertible elements of A(Φ). The algebra A(Φ) is a commutative algebra of degree 2 over A. When E contains no nonzero isotropic vector, A(Φ) is a field, a quadratic extension of A; otherwise it is the direct product of two fields isomorphic to A.
\item
  Let \((e_{1}, e_{2})\) an orthogonal basis of E; set \(\alpha_{i} = \Phi(e_{i}, e_{i})\) (i = 1, 2) and \(\delta = -\alpha_{2}/\alpha_{1}\) Then the matrices of the elements of A( \(\Phi\) ) with respect to this basis are the matrices of the form \(\begin{pmatrix} a & \delta b \\ b & a \end{pmatrix}\) where \(a \in A\) , \(b \in A\) .
\item
  The space E is a free cyclic A(Φ)-module, generated by any non-isotropic vector.
\end{enumerate}

Indeed, let us introduce an auxiliary alternating bilinear form B ≠ 0 on E; then B is non-degenerate. There therefore exists an endomorphism w of E such that \(\Phi(x, y) = \mathrm{B}(w(x), y)\) for all x, y in E. For every invertible endomorphism u of E, we have

\[
\begin{array}{r l} \Phi (u (x), u (y)) & = \mathrm{B} (w u (x), u (y)) = \mathrm{B} (u u ^ {- 1} w u (x), u (y)) \\ & = (\det u) \mathrm{B} (u ^ {- 1} w u (x), y); \end{array}
\]

For u to be a direct similarity, it is necessary and sufficient that one have

\[
(\det u) \mathrm{B} (u ^ {- 1} w u (x), y) = (\det u) \Phi (x, y) = (\det u) \mathrm{B} (w (x), y)
\]

for all \(x, y\) in E; since det \(u \neq 0\) and since B is nondegenerate, this is equivalent to \(u^{-1}wu = w\) , or equivalently to \(uw = wu\) . Take for B the alternating bilinear form whose matrix \(S\) with respect to \((e_1, e_2)\) is \(\begin{pmatrix} 0 & 1 \\ -1 & 0 \end{pmatrix}\) by denoting \(R\) and \(W\) the matrices of \(\Phi\) and of \(w\) with respect to this basis, the relation \(\Phi(x, y) = \mathrm{B}(w(x), y)\) is written, by virtue of formula (47) of § 1, no. 10, \(R = ^t W.S\) ; making this explicit shows that we have \(W = \begin{pmatrix} 0 & \alpha_2 \\ -\alpha_1 & 0 \end{pmatrix}\) . If \(\begin{pmatrix} a & c \\ b & d \end{pmatrix}\) denotes the matrix of \(u\) with respect to \((e_1, e_2)\) , the relation \(uw = wu\) is therefore equivalent to the relations \(b\alpha_2 = -c\alpha_1\) , \(a\alpha_2 = d\alpha_2\) , \(a\alpha_1 = d\alpha_1\) that is, to \(a = d\) and \(c = \delta b\) ; this proves that the matrices of direct similarities are the invertible matrices of the form \(\begin{pmatrix} a & \delta b \\ b & a \end{pmatrix}(a, b\) in A).

Now the endomorphisms of E whose matrices are of the form \(\begin{pmatrix} a & \delta b \\ b & a \end{pmatrix}\) (a, b in A) form a vector subspace of dimension 2 of \(\mathfrak{L}_{\mathrm{A}}(\mathrm{E})\) generated by 1 and by the endomorphism w; since \(w^{2}\) is the homothety with ratio \(-\alpha_{1}\alpha_{2}\) , this subspace is the subalgebra \(\mathrm{A}(\Phi)\) of \(\mathfrak{L}_{\mathrm{A}}(\mathrm{E})\) generated by the direct similarities. The direct similarities are the invertible elements of \(\mathrm{A}(\Phi)\) , that is, those whose matrices satisfy \(a^{2}-\delta b^{2}\neq0\) . The fact that the algebra \(\mathrm{A}(\Phi)\) is commutative is evident. Let us apply to it the results of Chap.~II, §7, no.~7: if δ is not a square in A, that is, if no nonzero vector of E is isotropic, \(\mathrm{A}(\Phi)\) is a field; if, on the contrary, δ is a square in A, that is, if E contains nonzero isotropic vectors, \(\mathrm{A}(\Phi)\) is a direct sum of two fields isomorphic to A. This proves a) and b).

Finally, any non-isotropic vector of E can be taken as

first vector \(e_1\) of an orthogonal basis \((e_1, e_2)\) ( \(\S 6, n^0 1\) ); therefore its transforms \(u(e_1)\) by the elements \(u\) of A( \(\Phi\) ) are the vectors of the form \(ae_1 + be_2\) ( \(a, b\) in A), that is, all vectors of E, since all matrices \(\begin{pmatrix} a & \delta b \\ b & a \end{pmatrix}\) are matrices of elements of A( \(\Phi\) ). In other words, E is a cyclic A( \(\Phi\) )-module, generated by any non-isotropic vector. Moreover, we see that it is a A( \(\Phi\) -module monogenic free, since \(u(e_1) = ae_1 + be_2 = 0\) implies \(a = b = 0\) , hence \(u = 0\) . This proves \(c\) ).

Remarks. --- 1) Let \(\vartheta\) be the similarity with matrix \(\begin{pmatrix}0 & \delta \\ 1 & 0\end{pmatrix}\) with respect to \((e_1, e_2)\) ; the similarity \(w\) introduced in the proof of Prop. 1 is equal to \(-\alpha_1 v\) ; we have \(v^2 = \delta\) . The multiplier of the direct similarity \(u = a + b v\) ( \(a, b\) in A) is equal to the determinant of its matrix \(\begin{pmatrix}a & \delta b \\ b & a\end{pmatrix}\) that is, to \(a^2 - \delta b^2 = (a + b v)(a - b v) = u.\bar{u}\) , denoting by \(\bar{u}\) the conjugate of \(u\) in the algebra A(Φ) (Chap.~II, § 7, No. 7); in other words, the multiplier of \(u\) is the norm N( \(u\) ) of \(u\) in the algebra A(Φ) (ibid.). In particular, for a direct similarity \(u\) to be a rotation, it is necessary and sufficient that N( \(u\) ) = 1; for \(u\) to be a homothety, it is necessary and sufficient that \(u \in A^*\) .

\begin{enumerate}
\def\labelenumi{\arabic{enumi})}
\setcounter{enumi}{1}
\item
  The direct similarities \(u = a + bv\) ( \(a, b\) in A, \(b \neq 0\) ) whose square is a homothety are the homotheties and the scalar multiples \(bv\) of \(v\) , since \((a + bv)^{2} = (a^{2} + \delta b^{2}) + 2abv\) . These latter are none other than the automorphisms of the vector space E which transform every vector into an orthogonal vector; indeed the matrix of such an automorphism is necessarily of the form \(\begin{pmatrix} 0 & c \\ d & 0 \end{pmatrix} (c, d\) in A), and the condition of transforming the vector \(\lambda e_{1} + \mu e_{2}\) as an orthogonal vector is then written \(\lambda \mu(d\alpha_{2} + c\alpha_{1}) = 0\) .
\item
  One easily verifies that, for \(x\) , \(y\) in E and \(u \in A(\Phi)\) , one has \(\Phi(u(x), y) =: \Phi(x, \overline{u}(y))\) . Thus the adjoint endomorphism of a direct similarity \(u\) is the direct similarity \(\overline{u}\) conjugate of \(u\) in A( \(\Phi\) ).
\item
  Since every inverse similarity of E is the product of a direct similarity and the reflection with respect to \(Ae_{1}\) , the matrices of inverse similarities with respect to \((e_{1}, e_{2})\) are the matrices of the form \(\begin{pmatrix} a & -\delta b \\ b & -a \end{pmatrix}\) .
\end{enumerate}

We shall henceforth denote by S the group of similarities of E, by \(S^{+}\) the group of direct similarities, by H that of homotheties ≠ 0, and by \(O^{+}\) that of rotations. Recall that we have \(H \subset S^{+}\) (§ 6, no. 5).

Corollary 1. --- The group S \(^{+}\) of direct similarities is commutative. For any non-isotropic vectors x, y in E, there exists a unique direct similarity u such that y = u(x).

The first assertion follows from the fact that the algebra A(Φ) is commutative. Since E is a free monogenic A(Φ)-module generated by x (resp. y), there exists a unique element u (resp. u′) of A(Φ) such that \(y = u(x)\) (resp. \(x = u'(y)\) ); whence \(x = u(u'(x))\) , and \(uu'\) is the identity; this shows that u is invertible, and is therefore a direct similarity.

COROLLARY 2. --- The group O \(^{+}\) of rotations is commutative. For any vectors x, y in E such that Φ(x, x) = Φ(y, y) ≠ 0, there exists a unique rotation u such that y = u(x).

The first assertion follows from Cor. 1. This also shows that there exists a direct similarity \(u\) and exactly one such that \(y = u(x)\) ; since \(\Phi(u(x), u(x)) = \Phi(x, x)\) the multiplier of \(u\) is equal to 1, and \(u\) is therefore a rotation.

COROLLARY 3. --- The group S \(^{+}\) /H is commutative. It acts on the set of non-isotropic lines of E. Whatever the non-isotropic lines D, D′ of E may be, there exists one and only one element φ of S \(^{+}\) /H such that D' = φ(D).

This follows from Cor. 1 and from the fact that H leaves every line of E globally invariant.

PROPOSITION 2. --- The kernel of the canonical homomorphism from O \(^{+}\) into S \(^{+}\) /H is \{1, -1\}.

Indeed, 1 and -1 are the only elements of \(\mathbf{H} \cap \mathbf{O}^{+}\) .

PROPOSITION 3. --- The homomorphism \(u \to u/\overline{u} = u^{2}/\mathrm{N}(u)\) of \(\mathrm{S}^{+}\) into itself has H as its kernel, and defines, by passing to the quotient, an isomorphism from \(\mathrm{S}^{+}/\mathrm{H}\) onto \(\mathrm{O}^{+}\) .

Indeed, the relation \(u/\overline{u} = 1\) is equivalent to \(u = \overline{u}\) that is, to \(u \in A^* = H\) Thus H is the kernel of \(u \to u/\overline{u}\) . Since \(N(u/\overline{u}) = 1\) , \(u/\overline{u}\) is a rotation (Remark 1). It remains to show that every element \(v\) of \(O^+\) is of the form \(u/\overline{u}\) ( \(u \in S^+\) ). If \(1 + v\) is invertible, one can take \(u = 1 + v\) , because the relation \(N(v) = v\bar{v} = 1\) implies \(1 + v = v(1 + \bar{v})\) . Otherwise, one has \(N(1 + v) = (1 + v)(1 + \bar{v}) = 0\) that is, setting \(v = a + bw\) ( \(a \in A\) , \(b \in A\) , \(w^2 = \delta \in A\) ), \(2(1 + a) = 0\) , whence \(a = -1\) ; but the relations \(a = -1\) and \(N(v) = a^2 - \delta b^2 = 1\) imply \(b = 0\) , whence \(v = -1\) ; since \(\bar{w} = -w\) , it suffices, in this case, to take \(u = w\) .

When A(Φ) is a field, prop. 3 is a special case of Hilbert's theorem (chap.~V, § 11, no. 5, th. 3).

COROLLARY. — Let \(i: \mathrm{O}^{+} \to \mathrm{S}^{+}/\mathrm{H}\) and \(d: \mathrm{S}^{+}/\mathrm{H} \to \mathrm{O}^{+}\) denote the homomorphisms defined in props. 2 and 3, and write the abelian groups \(\mathrm{O}^{+}\) and \(\mathrm{S}^{+}/\mathrm{H}\) additively; we have \(d(i(\theta)) = 2\theta\) for \(\theta \in \mathrm{O}^{+}\) and \(i(d(\varphi)) = 2\varphi\) for \(\varphi \in \mathrm{S}^{+}/\mathrm{H}\) .

Indeed, if \(\bar{\nu}\) is a rotation, we have \(\bar{\nu} = \nu^{-1}\) , whence \(d(i(\nu)) = \nu/\bar{\nu} = \nu^{2}\) . On the other hand, if \(\varphi \in S^{+}/H\) , \(\varphi\) is the class mod. H of a direct similarity \(u\) , and \(d(\varphi) = u/\bar{u} = u^{2}/N(u)\) is congruent to \(u^{2}\) mod. H; whence the second formula.

